\SourceLeaf{032}
circumferentiam in eo puncto inter (3) et (5), ubi $T=0$, revera \textit{secari} ab aliquo ramo lineae primae; idemque valet de reliquis punctis, ubi $T=0$. Eodem modo patet. circumferentiam circuli huius in omnibus $2m$ punctis, ubi $U=0$, ab aliquo ramo lineae secundae secari. Hae conclusiones etiam sequenti modo exprimi possunt: Descripto circulo debitae magnitudinis e centro $C$, in hunc intrabunt $2m$ rami lineae primae totidemque rami lineae secundae, et quidem ita, ut bini rami proximi lineae primae per aliquem ramum lineae secundae ab invicem separentur. Vid. fig. 2, ubi circulus iam non infinitae sed finitae magnitudinis erit, numerique singulis ramis adscripti cum numeris, per quos in art. praec. et hoc limites certos in peripheria brevitatis caussa designavi, non sunt confundendi.

\Article{21} Iam ex hoc situ relativo ramorum in circulum intrantium tot modis diversis deduci potest, intersectionem alicuius rami lineae primae cum ramo lineae secundae intra circulum necessario dari, ut, quaenam potissimum methodus prae reliquis eligenda sit, propemodum nesciam. Luculentissima videtur esse haec: Designemus (fig. 2) punctum peripheriae circuli, ubi a laeva axis parte (quae ipsa est unus ex $2m$ ramis lineae primae) secatur, per 0; punctum proximum, ubi ramus lineae secundae intrat, per 1; punctum huic proximum, ubi secundus lineae primae ramus intrat, per 2, et sic porro usque ad $4m-1$, ita ut in quovis puncto numero pari signato ramus lineae secundae in circulum intret, contra ramus lineae secundae in omnibus punctis per numerum imparem expressis. Iam ex geometria sublimiori constat, quamvis curvam algebraicam, (sive singulas cuiusvis curvae algebraicae partes, si forte e pluribus composita sit) aut in se redeuntem aut utrimque in infinitum excurrentem esse, adeoque si ramus aliquis curvae algebraicae in spatium definitum intret, eundem necessario ex hoc spatio rursus alicubi exire debere\HistoricalNote{10}. Hinc concluditur facile, quodvis punctum numero pari \SourceLeaf{033}
signatum (seu, brevitatis caussa, \textit{quodvis punctum par}) per ramum lineae primae cum alio puncto pari intra circulum iunctum esse debere, similiterque quodvis punctum numero impari notatum cum alio simili puncto per ramum lineae secundae. Quamquam vero haec binorum punctorum connexio secundum indolem functionis $X$ perquam diversa esse potest, ita ut in genere determinari nequeat, tamen facile demonstrari potest, \textit{quaecunque demum illa sit, semper intersectionem lineae primae cum linea secunda oriri.}

\Article{22} Demonstratio huius necessitatis commodissime apagogice repraesentari posse videtur. Scilicet supponamus, iunctionem binorum quorumque punctorum parium, et binorum quorumque punctorum imparium ita adornari posse, ut nulla intersectio rami lineae primae cum ramo lineae secundae inde oriatur. Quoniam axis est pars lineae primae, manifesto punctum 0 cum puncto $2m$ iunctum erit. Punctum 1 itaque cum nullo puncto ultra axem sito, i. e. cum nullo puncto per numerum maiorem quam $2m$ expresso iunctum esse potest, alioquin enim linea iungens necessario axem secaret. Si itaque 1 cum puncto $n$ iunctum esse supponitur, erit $n<2m$. Ex simili ratione si 2 cum $n'$ iunctum esse statuitur, erit $n'<n$, quia alioquin ramus $2\ldots n'$ ramum $1\ldots n$ necessario secaret. Ex eadem caussa punctum 3 cum aliquo punctorum inter 4 et $n'$ iacentium iunctum erit, patetque si 3, 4, 5 etc. iuncta esse supponantur cum $n''$, $n'''$, $n''''$ etc., $n'''$ iacere inter 5 et $n''$, $n''''$ inter 6 et $n'''$ etc. Unde perspicuum est, tandem ad aliquod punctum $h$ perventum iri, quod cum puncto $h+2$ iunctum sit, et tum ramus, qui in puncto $h+1$ in circulum intrat, necessario ramum puncta $h$ et $h+2$ iungentem secabit. Quia autem alter horum duorum ramorum ad lineam primam, alter ad secundam pertinebit, manifestum iam est, suppositionem esse contradictoriam, adeoque necessario alicubi intersectionem lineae primae cum linea secunda fieri.
\SourceLeaf{034}
Si haec cum praecedentibus iunguntur. ex omnibus disquisitionibus explicatis colligetur, theorema, \textit{quamvis functionem algebraicam rationalem integram unius indeterminatae in factores reales primi vel secundi gradus resolvi posse}, omni rigore demonstratum.

\Article{23} Ceterum haud difficile ex iisdem principiis deduci potest, non solum unam sed ad minimum $m$ intersectiones lineae primae cum secunda dari, quamquam etiam fieri potest, ut linea prima a pluribus ramis lineae secundae in eodem puncto secetur, in quo casu functio $X$ plures factores aequales habebit. Attamen quum hic sufficiat, unius intersectionis necessitatem demonstravisse, fusius huic rei brevitatis caussa non immoror. Ex eadem ratione etiam alias harum linearum proprietates hic uberius non persequor, e. g. intersectionem semper fieri sub angulis rectis; aut si plura crura utriusque curvae in eodem puncto conveniant, totidem crura lineae primae affore, quot crura lineae secundae, haecque alternatim posita esse, et sub aequalibus angulis se secare etc.

Denique observo, minime impossibile esse, ut demonstratio praecedens, quam hic principiis geometricis superstruxi, etiam in forma mere analytica exhibeatur: sed eam repraesentationem, quam hic explicavi, minus abstractam evadere credidi, verumque nervum probandi hic multo clarius ob oculos poni, quam a demonstratione analytica exspectari possit.

Coronidis loco adhuc aliam methodum theorema nostrum demonstrandi addigitabo, quae primo aspectu non modo a demonstratione praecedente, sed etiam ab omnibus demonstrationibus reliquis supra enarratis maxime diversa esse videbitur, et quae nihilominus cum \textsc{d'Alembert}iana, si ad essentiam spectas proprie eadem est. Cum qua illam comparare, parallelismumque inter utramque explorare peritis committo, in quorum gratiam unice subiuncta est.

\Article{24} Supra planum figurae 4 relative ad axem $CG$ punctumque fixum $C$ descriptas suppono superficiem primam et secundam eodem modo ut supra. Accipe punctum quodcunque in aliquo ramo lineae primae situm sive ubi $T=0$, (e. g. quodlibet punctum $M$ in axe iacens), et nisi in hoc etiam $U=0$, progredere ex hoc puncto in linea prima versus eam partem, versus quam magnitudo \SourceLeaf{035}
absoluta ipsius $U$ decrescit. Si forte in puncto $M$ valor absolutus ipsius $U$ versus utramque partem decrescit, arbitrarium est, quorsum progrediaris; quid vero faciendum sit, si $U$ versus utramque partem crescat, statim docebo. Manifestum est itaque, dum semper in linea prima progrediaris, necessario tandem te ad punctum perventurum, ubi $U=0$, aut ad tale, ubi valor ipsius $U$ fiat minimum, e. g. punctum $N$. In priori casu quod quaerebatur, inventum est; in posteriori vero demonstrari potest, in hoc puncto plures ramos lineae primae sese intersecare (et quidem multitudinem parem ramorum), quorum semissis ita comparati sint, ut si in aliquem eorum deflectas (sive huc sive illuc) valor ipsius $U$ adhucdum decrescere pergat. (Demonstrationem huius theorematis, prolixiorem quam difficiliorem brevitatis gratia supprimere debeo.) In hoc itaque ramo iterum progredi poteris, donec $U$ aut fiat $=0$ (uti in fig. 4 evenit in $P$), aut denuo minimum. Tum rursus deflectes, necessarioque tandem ad punctum pervenies, ubi sit $U=0$.

Contra hanc demonstrationem obiici posset dubium, annon possibile sit, ut quantumvis longe progrediaris, et quamvis valor ipsius $U$ semper decrescat, tamen haec decrementa continuo tardiora fiant, et nihilominus ille valor limitem aliquem nusquam attingat; quae obiectio responderet quartae in art. 6. Sed haud difficile foret, terminum aliquem assignare, quem simulac transieris, valor ipsius $U$ necessario non modo semper rapidius mutari debeat, sed etiam \textit{decrescere} non amplius possit, ita ut antequam ad hunc terminum perveneris, necessario valor 0 iam affuisse debeat. Hoc vero et reliqua, quae in hac demonstratione addigitare tantummodo potui, alia occasione fusius exsequi mihi reservo.

\begin{center}
\rule{8cm}{0.3pt}\\[8pt]
\textit{Principia quibus haecce demonstratio innititur deteximus Initio Octob.} 1797.\\[8pt]
\rule{5cm}{0.3pt}
\end{center}
